\documentclass[10pt]{article}
\usepackage[margin=2.2cm]{geometry}
\usepackage[T1]{fontenc}
\usepackage{newtxtext,newtxmath}\usepackage[varqu,scaled=0.95]{zi4}
\usepackage{enumitem}
\setlist{itemsep=2pt, topsep=3pt}
\ifdefined\anonbuild\newcommand{\tool}{QLine}\else\newcommand{\tool}{srcSlice-Q}\fi
% Generated by experiments/icpc/make_numbers.py from committed results. Do not edit by hand.

\newcommand{\numNPairs}{784}  % results/icpc/rq1/records.csv
\newcommand{\numNCircuits}{158}  % results/icpc/rq1/records.csv
\newcommand{\numRqOneDuUnsound}{216}  % results/icpc/rq1/records.csv
\newcommand{\numRqOneDuUnsoundPct}{27.6\%}  % results/icpc/rq1/records.csv
\newcommand{\numRqOneDuRecall}{0.718}  % results/icpc/rq1/records.csv
\newcommand{\numRqOneDuPrecision}{0.970}  % results/icpc/rq1/records.csv
\newcommand{\numRqOneDuSize}{51.2\%}  % results/icpc/rq1/records.csv
\newcommand{\numRqOneLcUnsound}{0}  % results/icpc/rq1/records.csv
\newcommand{\numRqOneLcUnsoundPct}{0.0\%}  % results/icpc/rq1/records.csv
\newcommand{\numRqOneLcRecall}{1.000}  % results/icpc/rq1/records.csv
\newcommand{\numRqOneLcPrecision}{0.923}  % results/icpc/rq1/records.csv
\newcommand{\numRqOneLcSize}{70.4\%}  % results/icpc/rq1/records.csv
\newcommand{\numRqOneYzUnsound}{0}  % results/icpc/rq1/records.csv
\newcommand{\numRqOneYzUnsoundPct}{0.0\%}  % results/icpc/rq1/records.csv
\newcommand{\numRqOneYzRecall}{1.000}  % results/icpc/rq1/records.csv
\newcommand{\numRqOneYzPrecision}{0.835}  % results/icpc/rq1/records.csv
\newcommand{\numRqOneYzSize}{78.6\%}  % results/icpc/rq1/records.csv
\newcommand{\numRqOneQUnsound}{0}  % results/icpc/rq1/records.csv
\newcommand{\numRqOneQUnsoundPct}{0.0\%}  % results/icpc/rq1/records.csv
\newcommand{\numRqOneQRecall}{0.807}  % results/icpc/rq1/records.csv
\newcommand{\numRqOneQPrecision}{0.972}  % results/icpc/rq1/records.csv
\newcommand{\numRqOneQSize}{54.0\%}  % results/icpc/rq1/records.csv
\newcommand{\numRqOneFourEdgeUnfixedUnsound}{5}  % results/icpc/rq1/records.csv
\newcommand{\numRqOneFourEdgeUnfixedUnsoundPct}{0.6\%}  % results/icpc/rq1/records.csv
\newcommand{\numRqOneFourEdgeUnfixedRecall}{0.994}  % results/icpc/rq1/records.csv
\newcommand{\numRqOneFourEdgeUnfixedPrecision}{0.864}  % results/icpc/rq1/records.csv
\newcommand{\numRqOneFourEdgeUnfixedSize}{75.0\%}  % results/icpc/rq1/records.csv
\newcommand{\numRqOneFourEdgeFixedUnsound}{0}  % results/icpc/rq1/records.csv
\newcommand{\numRqOneFourEdgeFixedUnsoundPct}{0.0\%}  % results/icpc/rq1/records.csv
\newcommand{\numRqOneFourEdgeFixedRecall}{1.000}  % results/icpc/rq1/records.csv
\newcommand{\numRqOneFourEdgeFixedPrecision}{0.864}  % results/icpc/rq1/records.csv
\newcommand{\numRqOneFourEdgeFixedSize}{75.2\%}  % results/icpc/rq1/records.csv
\newcommand{\numRqOneDuMissKickback}{1{,}737}  % results/icpc/rq1/misses.csv
\newcommand{\numRqOneDuMissSameWire}{835}  % results/icpc/rq1/misses.csv
\newcommand{\numRqOneDuMissWrittenOperand}{654}  % results/icpc/rq1/misses.csv
\newcommand{\numRqOneDuMissFeedback}{134}  % results/icpc/rq1/misses.csv
\newcommand{\numRqOneQRemoved}{2{,}847}  % results/icpc/rq1/pairs.csv
\newcommand{\numRqOneQRemovedNonInfluencing}{669}  % results/icpc/rq1/pairs.csv
\newcommand{\numRqOneQRemovedContributing}{1{,}618}  % results/icpc/rq1/pairs.csv
\newcommand{\numRqOneQPrunedPairs}{477}  % results/icpc/rq1/pairs.csv
\newcommand{\numLcqApplicable}{732}  % results/icpc/rq1_lcq_remap/records.csv
\newcommand{\numLcqNotApplicable}{52}  % results/icpc/rq1_lcq_remap/records.csv
\newcommand{\numLcqNaMeasurements}{42}  % results/icpc/rq1_lcq_remap/records.csv
\newcommand{\numLcqUnsound}{0}  % results/icpc/rq1_lcq_remap/records.csv
\newcommand{\numLcqSmallerThanLc}{243}  % results/icpc/rq1_lcq_remap/records.csv
\newcommand{\numLcqRecall}{0.93}  % results/icpc/rq1_lcq_remap/records.csv
\newcommand{\numQVsLcqReduction}{12.3\%}  % results/icpc/rq1_lcq_remap/records.csv
\newcommand{\numQVsLcqReductionPooled}{10.8\%}  % results/icpc/rq1_lcq_remap/records.csv
\newcommand{\numQLargerThanLcq}{72}  % results/icpc/rq1_lcq_remap/records.csv
\newcommand{\numHOnePHolm}{\ensuremath{p < 10^{-10}}}  % results/icpc/confirmatory/summary.json
\newcommand{\numHTwoDiff}{16.3}  % results/icpc/confirmatory/summary.json
\newcommand{\numHTwoCILow}{15.1}  % results/icpc/confirmatory/summary.json
\newcommand{\numHTwoCIHigh}{17.7}  % results/icpc/confirmatory/summary.json
\newcommand{\numHTwoRelative}{23.2\%}  % results/icpc/confirmatory/summary.json
\newcommand{\numHTwoSmaller}{477}  % results/icpc/confirmatory/summary.json
\newcommand{\numHTwoPHolm}{\ensuremath{p < 10^{-10}}}  % results/icpc/confirmatory/summary.json
\newcommand{\numHTwoR}{1.00}  % results/icpc/confirmatory/summary.json
\newcommand{\numHThreeRate}{55.4\%}  % results/icpc/confirmatory/summary.json
\newcommand{\numHThreeCILow}{51.9\%}  % results/icpc/confirmatory/summary.json
\newcommand{\numHThreeCIHigh}{58.8\%}  % results/icpc/confirmatory/summary.json
\newcommand{\numHThreePrunedRate}{26.6\%}  % results/icpc/confirmatory/summary.json
\newcommand{\numHThreePrunedRobust}{127}  % results/icpc/confirmatory/summary.json
\newcommand{\numHThreePrunedN}{477}  % results/icpc/confirmatory/summary.json
\newcommand{\numHFourContainQ}{82.2\%}  % results/icpc/confirmatory/summary.json
\newcommand{\numHFourContainLc}{95.7\%}  % results/icpc/confirmatory/summary.json
\newcommand{\numHFourDiff}{\ensuremath{-}13.5}  % results/icpc/confirmatory/summary.json
\newcommand{\numHFourCILow}{\ensuremath{-}14.4}  % results/icpc/confirmatory/summary.json
\newcommand{\numHFourCIHigh}{\ensuremath{-}12.6}  % results/icpc/confirmatory/summary.json
\newcommand{\numHFiveSizeDiff}{2.9}  % results/icpc/confirmatory/summary.json
\newcommand{\numHFiveExamDiff}{\ensuremath{-}0.061}  % results/icpc/confirmatory/summary.json
\newcommand{\numRobQPrunedBasis}{65.8\%}  % results/icpc/rq2_robustness/pairs.csv
\newcommand{\numRobLcBasis}{100.0\%}  % results/icpc/rq2_robustness/pairs.csv
\newcommand{\numRobQPrunedHaar}{26.6\%}  % results/icpc/rq2_robustness/pairs.csv
\newcommand{\numRobLcHaar}{100.0\%}  % results/icpc/rq2_robustness/pairs.csv
\newcommand{\numRobQPrunedParam}{78.1\%}  % results/icpc/rq2_robustness/pairs.csv
\newcommand{\numRobLcParam}{100.0\%}  % results/icpc/rq2_robustness/pairs.csv
\newcommand{\numRobLcqBasisRobust}{732}  % results/icpc/rq2_robustness_lcq/pairs.csv
\newcommand{\numRobLcqBasisN}{732}  % results/icpc/rq2_robustness_lcq/pairs.csv
\newcommand{\numRobLcqHaarRobust}{732}  % results/icpc/rq2_robustness_lcq/pairs.csv
\newcommand{\numRobLcqHaarN}{732}  % results/icpc/rq2_robustness_lcq/pairs.csv
\newcommand{\numRobLcqParamRobust}{171}  % results/icpc/rq2_robustness_lcq/pairs.csv
\newcommand{\numRobLcqParamN}{175}  % results/icpc/rq2_robustness_lcq/pairs.csv
\newcommand{\numAblUnfixedFull}{0.5\%}  % results/rq4_ablation/TAB_ablation_rerun.csv
\newcommand{\numAblUnfixedNoED}{54.0\%}  % results/rq4_ablation/TAB_ablation_rerun.csv
\newcommand{\numAblUnfixedNoMD}{13.4\%}  % results/rq4_ablation/TAB_ablation_rerun.csv
\newcommand{\numAblUnfixedNoUED}{13.3\%}  % results/rq4_ablation/TAB_ablation_rerun.csv
\newcommand{\numAblUnfixedNoCD}{4.8\%}  % results/rq4_ablation/TAB_ablation_rerun.csv
\newcommand{\numAblFixedFull}{0.0\%}  % results/icpc/ablation_fixed/TAB_ablation_rerun.csv
\newcommand{\numAblFixedNoED}{63.9\%}  % results/icpc/ablation_fixed/TAB_ablation_rerun.csv
\newcommand{\numAblFixedNoMD}{13.0\%}  % results/icpc/ablation_fixed/TAB_ablation_rerun.csv
\newcommand{\numAblFixedNoUED}{26.1\%}  % results/icpc/ablation_fixed/TAB_ablation_rerun.csv
\newcommand{\numAblFixedNoCD}{6.7\%}  % results/icpc/ablation_fixed/TAB_ablation_rerun.csv
\newcommand{\numAblPairs}{962}  % results/icpc/ablation_fixed/TAB_ablation_rerun.csv
\newcommand{\numFamCircuits}{119}  % results/icpc/families/rows.csv
\newcommand{\numFamLcUnsound}{0}  % results/icpc/families/rows.csv
\newcommand{\numFamYzUnsound}{0}  % results/icpc/families/rows.csv
\newcommand{\numFamQUnsound}{0}  % results/icpc/families/rows.csv
\newcommand{\numFamDuUnsound}{232}  % results/icpc/families/rows.csv
\newcommand{\numFamQSoundNotMinimal}{212}  % results/icpc/families/rows.csv
\newcommand{\numFamQRows}{819}  % results/icpc/families/rows.csv
\newcommand{\numCatEntries}{24}  % results/icpc/catalogue/catalogue.json
\newcommand{\numRqThreeUnits}{5{,}688}  % results/icpc/rq3/records.csv
\newcommand{\numRqThreeContainDu}{77.3\%}  % results/icpc/rq3/records.csv
\newcommand{\numRqThreeContainLc}{95.7\%}  % results/icpc/rq3/records.csv
\newcommand{\numRqThreeContainYz}{97.4\%}  % results/icpc/rq3/records.csv
\newcommand{\numRqThreeContainQ}{82.2\%}  % results/icpc/rq3/records.csv
\newcommand{\numRqThreeContainRandom}{62.4\%}  % results/icpc/rq3/records.csv
\newcommand{\numCQFour}{769}  % results/icpc/rq3/units.csv
\newcommand{\numCQFourRecovered}{691}  % results/icpc/rq3/units.csv
\newcommand{\numBugsFourQKilled}{7}  % results/icpc/rq3_bugs4q/bugs.csv
\newcommand{\numBugsFourQInjured}{1}  % results/icpc/rq3_bugs4q/bugs.csv
\newcommand{\numBugsFourQAlive}{1}  % results/icpc/rq3_bugs4q/bugs.csv
\newcommand{\numBugsFourQExcluded}{2}  % results/icpc/rq3_bugs4q/bugs.csv
\newcommand{\numBugsFourQTriaged}{107}  % bugs4q/triage.csv
\newcommand{\numBugsFourQIncluded}{11}  % bugs4q/triage.csv
\newcommand{\numMapInstructions}{137}  % results/icpc/source_map/bugs.csv
\newcommand{\numMapMapped}{136}  % results/icpc/source_map/bugs.csv
\newcommand{\numMapCrossCompared}{127}  % results/icpc/source_map/crosscheck.csv
\newcommand{\numMapCrossAgree}{127}  % results/icpc/source_map/crosscheck.csv
\newcommand{\numMapCrossPrograms}{5}  % results/icpc/source_map/crosscheck.csv
\newcommand{\numLinesFullContained}{15}  % results/icpc/source_map/units.csv
\newcommand{\numLinesFullMean}{10.7}  % results/icpc/source_map/units.csv
\newcommand{\numLinesLcContained}{15}  % results/icpc/source_map/units.csv
\newcommand{\numLinesLcMean}{10.3}  % results/icpc/source_map/units.csv
\newcommand{\numLinesQContained}{15}  % results/icpc/source_map/units.csv
\newcommand{\numLinesQMean}{9.5}  % results/icpc/source_map/units.csv
\newcommand{\numLinesDuContained}{14}  % results/icpc/source_map/units.csv
\newcommand{\numLinesDuMean}{4.9}  % results/icpc/source_map/units.csv
\newcommand{\numLinesUnits}{15}  % results/icpc/source_map/units.csv
\newcommand{\numMbflUnits}{293}  % results/icpc/mbfl/units.csv
\newcommand{\numMbflExamAll}{0.348}  % results/icpc/mbfl/units.csv
\newcommand{\numMbflExamQ}{0.465}  % results/icpc/mbfl/units.csv
\newcommand{\numMbflExamLc}{0.419}  % results/icpc/mbfl/units.csv
\newcommand{\numTTwoSonnetFull}{84.4\%}  % results/llm/t2_scores.csv
\newcommand{\numTTwoSonnetLc}{77.2\%}  % results/llm/t2_scores.csv
\newcommand{\numTTwoSonnetQ}{69.9\%}  % results/llm/t2_scores.csv
\newcommand{\numTTwoSonnetDu}{62.3\%}  % results/llm/t2_scores.csv
\newcommand{\numTTwoSonnetQVsFullDiff}{\ensuremath{-}14.5}  % results/llm/t2_scores.csv
\newcommand{\numTTwoSonnetQVsFullCI}{[\ensuremath{-}19.9, \ensuremath{-}9.4]}  % results/llm/t2_scores.csv
\newcommand{\numTTwoSonnetQVsFullP}{\ensuremath{p = 1.6\times 10^{-7}}}  % results/llm/t2_scores.csv
\newcommand{\numTTwoSonnetQVsLcDiff}{\ensuremath{-}7.2}  % results/llm/t2_scores.csv
\newcommand{\numTTwoSonnetQVsLcCI}{[\ensuremath{-}12.3, \ensuremath{-}2.5]}  % results/llm/t2_scores.csv
\newcommand{\numTTwoSonnetQVsLcP}{\ensuremath{p = 0.006}}  % results/llm/t2_scores.csv
\newcommand{\numTTwoHaikuFull}{71.4\%}  % results/llm/t2_scores.csv
\newcommand{\numTTwoHaikuLc}{65.9\%}  % results/llm/t2_scores.csv
\newcommand{\numTTwoHaikuQ}{58.7\%}  % results/llm/t2_scores.csv
\newcommand{\numTTwoHaikuDu}{51.8\%}  % results/llm/t2_scores.csv
\newcommand{\numTTwoHaikuQVsFullDiff}{\ensuremath{-}12.7}  % results/llm/t2_scores.csv
\newcommand{\numTTwoHaikuQVsFullCI}{[\ensuremath{-}18.8, \ensuremath{-}6.5]}  % results/llm/t2_scores.csv
\newcommand{\numTTwoHaikuQVsFullP}{\ensuremath{p = 8.2\times 10^{-5}}}  % results/llm/t2_scores.csv
\newcommand{\numTTwoHaikuQVsLcDiff}{\ensuremath{-}7.2}  % results/llm/t2_scores.csv
\newcommand{\numTTwoHaikuQVsLcCI}{[\ensuremath{-}12.7, \ensuremath{-}2.2]}  % results/llm/t2_scores.csv
\newcommand{\numTTwoHaikuQVsLcP}{\ensuremath{p = 0.010}}  % results/llm/t2_scores.csv
\newcommand{\numTTwoQwenMidFull}{42.8\%}  % results/llm/t2_scores.csv
\newcommand{\numTTwoQwenMidLc}{43.1\%}  % results/llm/t2_scores.csv
\newcommand{\numTTwoQwenMidQ}{36.2\%}  % results/llm/t2_scores.csv
\newcommand{\numTTwoQwenMidDu}{33.3\%}  % results/llm/t2_scores.csv
\newcommand{\numTTwoQwenMidQVsFullDiff}{\ensuremath{-}6.5}  % results/llm/t2_scores.csv
\newcommand{\numTTwoQwenMidQVsFullCI}{[\ensuremath{-}11.6, \ensuremath{-}1.4]}  % results/llm/t2_scores.csv
\newcommand{\numTTwoQwenMidQVsFullP}{\ensuremath{p = 0.020}}  % results/llm/t2_scores.csv
\newcommand{\numTTwoQwenMidQVsLcDiff}{\ensuremath{-}6.9}  % results/llm/t2_scores.csv
\newcommand{\numTTwoQwenMidQVsLcCI}{[\ensuremath{-}11.2, \ensuremath{-}2.5]}  % results/llm/t2_scores.csv
\newcommand{\numTTwoQwenMidQVsLcP}{\ensuremath{p = 0.003}}  % results/llm/t2_scores.csv
\newcommand{\numTTwoQwenSmallFull}{13.0\%}  % results/llm/t2_scores.csv
\newcommand{\numTTwoQwenSmallLc}{12.7\%}  % results/llm/t2_scores.csv
\newcommand{\numTTwoQwenSmallQ}{13.0\%}  % results/llm/t2_scores.csv
\newcommand{\numTTwoQwenSmallDu}{13.4\%}  % results/llm/t2_scores.csv
\newcommand{\numTTwoQwenSmallQVsFullDiff}{0.0}  % results/llm/t2_scores.csv
\newcommand{\numTTwoQwenSmallQVsFullCI}{[\ensuremath{-}4.3, 4.3]}  % results/llm/t2_scores.csv
\newcommand{\numTTwoQwenSmallQVsFullP}{\ensuremath{p = 1.00}}  % results/llm/t2_scores.csv
\newcommand{\numTTwoQwenSmallQVsLcDiff}{0.4}  % results/llm/t2_scores.csv
\newcommand{\numTTwoQwenSmallQVsLcCI}{[\ensuremath{-}2.9, 3.6]}  % results/llm/t2_scores.csv
\newcommand{\numTTwoQwenSmallQVsLcP}{\ensuremath{p = 1.00}}  % results/llm/t2_scores.csv
\newcommand{\numTTwoItems}{276}  % results/llm/t2_scores.csv
\newcommand{\numCostLcMaxMs}{16}  % results/icpc/cost/circuits.csv
\newcommand{\numCostQDepthTen}{1.4}  % results/icpc/cost/circuits.csv
\newcommand{\numCostBuildDepthTen}{0.03}  % results/icpc/cost/circuits.csv
\newcommand{\numCostQDepthTwenty}{10.6}  % results/icpc/cost/circuits.csv
\newcommand{\numCostBuildDepthTwenty}{0.27}  % results/icpc/cost/circuits.csv
\newcommand{\numCostQDepthForty}{49.9}  % results/icpc/cost/circuits.csv
\newcommand{\numCostBuildDepthForty}{1.10}  % results/icpc/cost/circuits.csv
\newcommand{\numCostDepthEightyOverLimit}{4}  % results/icpc/cost/circuits.csv
\newcommand{\numCostBuildDepthEighty}{5.2}  % results/icpc/cost/circuits.csv
\newcommand{\numCostDepthOneSixtyOverLimit}{5}  % results/icpc/cost/circuits.csv
\newcommand{\numCostBuildDepthOneSixty}{21.7}  % results/icpc/cost/circuits.csv
\newcommand{\numCostWallLimit}{120}  % icpc/cost_run.py
\newcommand{\numCostWidthQMax}{8.7}  % results/icpc/cost/circuits.csv
\newcommand{\numCorpusLCircuits}{55}  % results/icpc/corpus_l/circuits.csv
\newcommand{\numCorpusLCompleted}{39}  % results/icpc/corpus_l/circuits.csv
\newcommand{\numNoiseLcCrossExactMax}{\ensuremath{4.2\times 10^{-15}}}  % results/icpc/noise/pairs.csv
\newcommand{\numNoiseLcCrossExactNonzero}{0}  % results/icpc/noise/pairs.csv
\newcommand{\numNoiseLcCrossExactN}{1{,}746}  % results/icpc/noise/pairs.csv
\newcommand{\numNoiseLcCrossShotMax}{0.033}  % results/icpc/noise/pairs.csv
\newcommand{\numNoiseQCrossExactMax}{\ensuremath{2.4\times 10^{-1}}}  % results/icpc/noise/pairs.csv
\newcommand{\numNoiseQCrossExactNonzero}{702}  % results/icpc/noise/pairs.csv
\newcommand{\numNoiseQCrossExactN}{1{,}746}  % results/icpc/noise/pairs.csv
\newcommand{\numNoiseQCrossShotMax}{0.033}  % results/icpc/noise/pairs.csv
\newcommand{\numNoiseShots}{1{,}024}  % icpc/noise_run.py
\newcommand{\numHwCircuits}{18}  % results/icpc/hardware/hardware_tvd_corrected.csv
\newcommand{\numHwWithin}{16}  % results/icpc/hardware/hardware_tvd_corrected.csv
\newcommand{\numHwNoiseFreeZero}{18}  % results/icpc/hardware/hardware_tvd_corrected.csv

\title{The Story of Paper B\\[2pt]\large Static Program Slicing for Quantum Programs}
\date{Big-picture document, 30 October 2026}
\begin{document}
\maketitle

\section*{What the paper is about}
Quantum programs are written as Python programs that build circuits. This paper presents an approach that slices
such programs on their source, without running them, and shows a developer the lines, the loop iterations and the
classical statements that determine a qubit. The approach is implemented in a tool, \tool{}, for Qiskit and
OpenQASM~2, and evaluated for faithfulness, coverage, cost and what it shows a developer.

\section*{The problem}
A Qiskit program decides in Python which gates act on which qubits: indices are computed, registers are sized by
variables, gates are added in loops and helper functions, and measured bits steer later operations. The tools a
developer has today see only part of this.
\begin{itemize}
\item A classical slicer reads the circuit as one object. To srcSlice, \texttt{qc.cx(a, b)} uses the variable
\texttt{qc}, so the slice of any qubit is every line that touches the circuit.
\item Circuit tools, such as Qiskit's DAG and its \texttt{LightCone} pass, work on a circuit that already exists. To get
one, the program must run, which needs its Qiskit version, its inputs and sometimes hardware credentials, and it
loses the link from each gate to the line and iteration that added it.
\item Analyses of quantum source code exist for OpenQASM~3, the srcML work of Behler et al., which names a qubit
slicer as future work. Qiskit, the language most quantum programs are written in, has bug detectors but no slicer.
\end{itemize}

\paragraph{Thesis.} Quantum programs can be sliced statically, on the source developers write, faithfully and at
scale. A semantic layer over srcML gives calls their quantum meaning, a static trace with over-approximations that
can only enlarge a slice stands in for running the program, and the slice points a developer to about a quarter of
a program.

\section*{Contributions}
\begin{enumerate}
\item \textbf{The approach.} A quantum layer over srcML and srcSlice with five parts: quantum roles of calls from a
table derived from Qiskit's interface (\numToolGateTable{} methods); a static trace of the circuit with loop iterations
and source lines, over-approximating unknown indices, branches and loop bounds; slices of every qubit along the
backward light cone extended through measured bits; the classical statements that compute each slice, from srcSlice
with two fixes; and a profile of each qubit.
\item \textbf{Loops that run on the device.} Qiskit's \texttt{while\_loop}, \texttt{for\_loop} and \texttt{switch} are
sliced by repeating the body, conditioned on its bits, until the needed set stops growing, with a soundness argument.
\item \textbf{Two front ends.} Qiskit programs through srcML, and OpenQASM~2 circuits through a line-mapped reader.
\item \textbf{An evaluation on real programs.} Faithfulness: on \numToolDevPrograms{} programs the trace equals the
circuit Qiskit builds and the slices equal circuit-level slices in \numToolDevSliceEqual{}; the OpenQASM reader agrees
with Qiskit on \numQasmCircuitsChecked{} circuits; exact simulation confirms every slice of \numCfQubits{} qubits in
\numCfPrograms{} programs with device loops; the classical statements cover a reference for \numSsContainPct{} of
qubits and \numSsHeldContainPct{} on held-out programs. Coverage: \numScaleSOneSliced{} of \numScaleSOneFiles{} GitHub
programs traced, \numCovNoConstruction{} of the rest never build a circuit; \numSubjBugsTraced{} of
\numSubjBugsCount{} Bugs4Q programs, \numSubjTextTraced{} of \numSubjTextCount{} textbook programs, and
\numQheTracedOfBuildsPct{} of the Qiskit HumanEval programs that build a circuit. Cost: \numScaleSOneMedianMs{}~ms per
program in the median; every qubit of a circuit of \numScaleQbMaxOps{} operations in \numScaleQbAllS{}~s, where Qiskit's
\texttt{LightCone} takes \numScaleLcMaxS{}~s for one qubit and the QPDG exceeds the limit beyond
\numScaleQpdgMaxOps{} operations. Value: a qubit's slice holds \numScaleReadMedianPct{} of a program's code lines; on
LintQ's warnings the slices give \numLqReadMedian{} of \numLqCodeMedian{} lines in the median, and
\numLqBeyondTenFiles{} traced programs create operations beyond the tenth loop iteration that LintQ unrolls.
\end{enumerate}

\section*{What is new, and what we build on}
\textbf{New.} The quantum layer for Qiskit source, its over-approximations and their soundness, slicing of loops that
run on the device, qubit profiles, and the first evaluation of a quantum slicer on thousands of real programs.

\textbf{Built on.} srcML's markup and srcSlice's slice profiles, patched for Python with two fixes; the backward light
cone, stated and cited as the dependence relation, not re-proved; Behler et al.'s idea of a static execution order of
qubit operations; Qiskit's own interface for the role table. Data from LintQ, Bugs4Q, the Qiskit textbook, Qiskit
HumanEval and QASMBench.

\section*{The motivating example}
A teleportation chain of four hops (\texttt{teleport\_chain.py}). The loop computes the qubits of each hop from its
counter, a variable \texttt{src} is reassigned on every iteration, and each hop measures two qubits whose bits steer
corrections on the next. A developer asks which lines determine the final qubit. srcSlice returns every line that
touches \texttt{qc}; \tool{} returns the lines, the iterations that put them there, and the statements that compute
\texttt{src}.

\section*{How the sections fit together}
I Introduction. II the motivating example. III Background and related work: srcML and srcSlice, analyses of quantum
source, circuit tools, bug detectors, and a table of approaches. IV The approach, with the overview figure and the
five parts, device loops and the two front ends. V Evaluation design: four research questions on faithfulness,
coverage, cost and value, the subjects, and the oracles. VI Results. VII Discussion of the limits, such as calls the
role table does not model. VIII Threats. IX Conclusion.

\section*{Rules for keeping Paper B separate}
It stands alone and does not cite the other submission. No sentence is reused from it; the material moved from it is
raw data, rewritten here. It has its own anonymous replication package. In the anonymous version the tool is
\tool{}.

\section*{Decisions to check}
\begin{enumerate}[label=(\roman*)]
\item \textbf{An independent faithfulness oracle.} The slices are now compared with a circuit-level slicer that the other
paper also uses. Comparing them with Qiskit's DAG ancestors instead would make the check independent; about three hours,
no paid runs.
\item \textbf{Semantic layer or quantum markup.} The approach interprets srcML's Python markup rather than emitting new
quantum elements. Describe it honestly as a semantic layer, or add emitted markup (16 to 25 hours).
\item \textbf{The example}, the teleportation chain above, or another program.
\item \textbf{Section order}, the example before related work, as in the other paper.
\item \textbf{Length.} The material is about 9 to 10 pages; the profile and LintQ comparison are the first to shorten.
\end{enumerate}
\end{document}
